\documentclass[11pt,a4paper]{article}
\usepackage[a4paper,margin=1in,headsep=30pt,footskip=35pt]{geometry}
\usepackage[T1]{fontenc}
\usepackage[utf8]{inputenc}
\usepackage{tgpagella}
\usepackage[scale=0.92]{tgheros}
\usepackage{microtype}
\usepackage{graphicx}
\usepackage{hyperref}
\usepackage{enumitem}
\usepackage{xcolor}
\usepackage{titlesec}
\usepackage{parskip}
\usepackage{longtable}
\usepackage{booktabs}
\usepackage{array}
\usepackage{tabularx}
\usepackage{fancyhdr}
\usepackage{lastpage}
\usepackage{tikz}
\usepackage[most]{tcolorbox}

% --- CORE COLOR IDENTITY ---
\definecolor{KazePrimary}{HTML}{284F58}
\definecolor{KazeSecondary}{HTML}{8A6538}
\definecolor{KazeInk}{HTML}{1F292D}
\definecolor{KazeMuted}{HTML}{5D676C}
\definecolor{KazeSoft}{HTML}{F8F4ED}
\definecolor{KazeLine}{HTML}{D8CCBC}
\definecolor{KazePaper}{HTML}{FCFBF8}
\definecolor{KazeLegal}{HTML}{7A2E2E}

\hypersetup{colorlinks=true, linkcolor=KazePrimary, urlcolor=KazePrimary}

\setlist[itemize]{leftmargin=1.5em, itemsep=0.4em}
\renewcommand{\arraystretch}{1.25}
\setlength{\headheight}{35pt}

\newcommand{\KazeLogoPath}{../../composeApp/src/commonMain/composeResources/drawable/k_logo_raster.png}
\newcommand{\KazeMarkPath}{../../composeApp/src/commonMain/composeResources/drawable/k_mark_raster.png}
\newcommand{\GaboLogoPath}{../../composeApp/src/commonMain/composeResources/drawable/gabo_mark_raster.png}
\newcommand{\GaboGrayMarkPath}{../../composeApp/src/commonMain/composeResources/drawable/gabo_mark_gray_raster.png}
\newcommand{\KazeWordmarkHeader}{{\fontfamily{qhv}\selectfont\sffamily\bfseries\small\textcolor{KazePrimary}{KAZE}}}
\newcommand{\KazeWordmarkCover}{{\fontfamily{qhv}\selectfont\sffamily\bfseries\fontsize{30}{30}\selectfont\textcolor{KazePrimary}{KAZE}}}
\newcommand{\KazeByGaboHeader}{{\raisebox{0.05cm}{\includegraphics[height=0.18cm]{\GaboLogoPath}}\hspace{0.2em}\sffamily\tiny\textcolor{KazeSecondary}{by GABO}}}
\newcommand{\KazeByGaboCover}{{\includegraphics[height=0.45cm]{\GaboLogoPath}\hspace{0.4em}\sffamily\large\textcolor{KazeSecondary}{by GABO}}}
\newcommand{\KazeCoverLabel}[1]{{\sffamily\normalsize\textcolor{KazeSecondary}{\textbf{#1}}}}
\newcommand{\KazeCoverTitle}[1]{{\fontfamily{qpl}\selectfont\fontsize{38}{43}\selectfont\textcolor{KazeInk}{#1}}}
\newcommand{\KazeCoverSubtitle}[1]{{\sffamily\Large\textcolor{KazePrimary}{#1}}}
\newcommand{\KazeContactEmail}{orestegabo@icloud.com}
\newcommand{\KazeContactWebsiteUrl}{https://kazerwanda.com}
\newcommand{\KazeContactWebsiteLabel}{kazerwanda.com}
\newcommand{\KazeContactBlock}{
    \begin{itemize}
        \item \textbf{Email:} \href{mailto:\KazeContactEmail}{\KazeContactEmail}
        \item \textbf{Website:} \href{\KazeContactWebsiteUrl}{\KazeContactWebsiteLabel}
    \end{itemize}
}

\pagecolor{KazePaper}

\AddToHook{shipout/background}{%
    \begin{tikzpicture}[remember picture,overlay]
        \fill[KazePrimary, opacity=0.018, rounded corners=14pt]
            ([xshift=0.04\paperwidth,yshift=-0.07\paperheight]current page.north west)
            rectangle
            ([xshift=0.96\paperwidth,yshift=-0.31\paperheight]current page.north west);
        \fill[KazeSecondary, opacity=0.014, rounded corners=14pt]
            ([xshift=0.05\paperwidth,yshift=0.34\paperheight]current page.north west)
            rectangle
            ([xshift=0.95\paperwidth,yshift=0.14\paperheight]current page.south west);

        \draw[KazePrimary, line width=1.7pt, opacity=0.09, line cap=round]
            ([xshift=0.08\paperwidth,yshift=-0.10\paperheight]current page.north west) --
            ([xshift=0.72\paperwidth,yshift=-0.10\paperheight]current page.north west);
        \draw[KazePrimary, line width=0.95pt, opacity=0.085, line cap=round]
            ([xshift=0.12\paperwidth,yshift=-0.135\paperheight]current page.north west) --
            ([xshift=0.88\paperwidth,yshift=-0.135\paperheight]current page.north west);
        \draw[KazeSecondary, line width=1.2pt, opacity=0.08, line cap=round]
            ([xshift=0.18\paperwidth,yshift=0.26\paperheight]current page.south west) --
            ([xshift=0.84\paperwidth,yshift=0.26\paperheight]current page.south west);
        \draw[KazePrimary, line width=1.7pt, opacity=0.09, line cap=round]
            ([xshift=0.14\paperwidth,yshift=0.21\paperheight]current page.south west) --
            ([xshift=0.62\paperwidth,yshift=0.21\paperheight]current page.south west);

        \draw[KazePrimary, line width=1.15pt, opacity=0.09]
            ([xshift=0.62\paperwidth,yshift=-0.06\paperheight]current page.north west) --
            ([xshift=0.76\paperwidth,yshift=-0.06\paperheight]current page.north west) --
            ([xshift=0.82\paperwidth,yshift=-0.11\paperheight]current page.north west) --
            ([xshift=0.93\paperwidth,yshift=-0.11\paperheight]current page.north west) --
            ([xshift=0.93\paperwidth,yshift=-0.18\paperheight]current page.north west) --
            ([xshift=0.82\paperwidth,yshift=-0.18\paperheight]current page.north west) --
            ([xshift=0.75\paperwidth,yshift=-0.24\paperheight]current page.north west) --
            ([xshift=0.58\paperwidth,yshift=-0.24\paperheight]current page.north west);

        \draw[KazeSecondary, line width=1.15pt, opacity=0.08]
            ([xshift=0.10\paperwidth,yshift=0.12\paperheight]current page.south west) --
            ([xshift=0.26\paperwidth,yshift=0.12\paperheight]current page.south west) --
            ([xshift=0.33\paperwidth,yshift=0.17\paperheight]current page.south west) --
            ([xshift=0.48\paperwidth,yshift=0.17\paperheight]current page.south west) --
            ([xshift=0.48\paperwidth,yshift=0.10\paperheight]current page.south west) --
            ([xshift=0.34\paperwidth,yshift=0.10\paperheight]current page.south west) --
            ([xshift=0.27\paperwidth,yshift=0.05\paperheight]current page.south west) --
            ([xshift=0.12\paperwidth,yshift=0.05\paperheight]current page.south west);

        \draw[KazePrimary, line width=1.8pt, opacity=0.05]
            ([xshift=0.70\paperwidth,yshift=-0.22\paperheight]current page.north west) circle [radius=2.0cm];
        \draw[KazeSecondary, line width=1.8pt, opacity=0.048]
            ([xshift=0.14\paperwidth,yshift=0.18\paperheight]current page.south west) circle [radius=2.35cm];

        \draw[KazePrimary, line cap=round, opacity=0.065, line width=1.2pt]
            ([xshift=0.11\paperwidth,yshift=-0.185\paperheight]current page.north west) --
            ([xshift=0.36\paperwidth,yshift=-0.185\paperheight]current page.north west);
        \draw[KazeSecondary, line cap=round, opacity=0.06, line width=0.8pt]
            ([xshift=0.11\paperwidth,yshift=-0.203\paperheight]current page.north west) --
            ([xshift=0.42\paperwidth,yshift=-0.203\paperheight]current page.north west);
        \draw[KazePrimary, line cap=round, opacity=0.065, line width=1.2pt]
            ([xshift=0.11\paperwidth,yshift=-0.221\paperheight]current page.north west) --
            ([xshift=0.48\paperwidth,yshift=-0.221\paperheight]current page.north west);
        \draw[KazeSecondary, line cap=round, opacity=0.06, line width=0.8pt]
            ([xshift=0.11\paperwidth,yshift=-0.239\paperheight]current page.north west) --
            ([xshift=0.54\paperwidth,yshift=-0.239\paperheight]current page.north west);
        \draw[KazePrimary, line cap=round, opacity=0.065, line width=1.2pt]
            ([xshift=0.11\paperwidth,yshift=-0.257\paperheight]current page.north west) --
            ([xshift=0.60\paperwidth,yshift=-0.257\paperheight]current page.north west);
        \draw[KazeSecondary, line cap=round, opacity=0.06, line width=0.8pt]
            ([xshift=0.11\paperwidth,yshift=-0.275\paperheight]current page.north west) --
            ([xshift=0.66\paperwidth,yshift=-0.275\paperheight]current page.north west);
        \draw[KazePrimary, line cap=round, opacity=0.065, line width=1.2pt]
            ([xshift=0.11\paperwidth,yshift=-0.293\paperheight]current page.north west) --
            ([xshift=0.72\paperwidth,yshift=-0.293\paperheight]current page.north west);
        \node[anchor=south east, opacity=0.035] at ([xshift=-0.45cm,yshift=0.65cm]current page.south east) {%
            \includegraphics[width=4.8cm]{\KazeMarkPath}%
        };
        \node[anchor=north east, opacity=0.028] at ([xshift=-1.1cm,yshift=-1.2cm]current page.north east) {%
            \includegraphics[width=2.2cm]{\GaboGrayMarkPath}%
        };
    \end{tikzpicture}%
}

% --- SECTION DESIGN: TECHNICAL PRECISION ---
\titleformat{\section}
{\sffamily\Large\bfseries\color{KazeInk}}
{\thesection}
{1em}
{\MakeUppercase}
[\vspace{0.2em}{\color{KazeSecondary}\hrule height 1.5pt width 2cm}\vspace{0.5em}]

\titleformat{\subsection}
{\sffamily\large\bfseries\color{KazePrimary}}
{\thesubsection}
{1em}
{}

% --- DUAL-BRANDED HEADER (HIGH VISIBILITY) ---
\pagestyle{fancy}
\fancyhf{}
\renewcommand{\headrulewidth}{0pt}
\renewcommand{\footrulewidth}{0pt}

\fancyhead[L]{%
    \begin{tabular}{@{}l l@{}}
        \raisebox{-0.3\height}{\includegraphics[height=0.6cm]{\KazeLogoPath}} &
        \begin{minipage}[b]{4cm}
            \KazeWordmarkHeader\\
            \vspace{-1mm}
            \KazeByGaboHeader
        \end{minipage}
    \end{tabular}
}
\fancyhead[R]{\sffamily\tiny\textcolor{KazeMuted}{KAZE PROPRIETARY // \today}}

\fancyfoot[L]{%
    \begin{tabular}{@{}l@{}}
        {\color{KazeLine}\rule{3.1cm}{0.6pt}}\\[-0.2em]
        {\sffamily\tiny\textcolor{KazeLegal}{CONFIDENTIAL PRIVATE DOCUMENT}}
    \end{tabular}
}
\fancyfoot[C]{%
    \sffamily\scriptsize
    \textcolor{KazeMuted}{\textsc{Page} \thepage\ of \pageref*{LastPage}}
}
\fancyfoot[R]{%
    \begin{tabular}{@{}r@{}}
        {\color{KazeLine}\rule{2.1cm}{0.6pt}}\\[-0.2em]
        {\sffamily\tiny\textcolor{KazePrimary}{Kaze by GABO}}
    \end{tabular}
}

% --- CLEAN REGISTRY BOXES ---
\newtcolorbox{kazehighlightbox}{
    enhanced,
    colback=KazeSoft,      % Only a very faint paper-like tint
    frame hidden,          % No borders at all
    sharp corners,
    left=5mm, right=5mm, top=4mm, bottom=4mm,
% Use a small square bullet as an anchor instead of a line
    overlay={
        \fill[KazeSecondary] ([xshift=2mm,yshift=-4mm]frame.north west) rectangle ([xshift=3mm,yshift=-5mm]frame.north west);
    }
}
\newtcolorbox{kazecalloutbox}[1][]{
    enhanced,
    colback=white,
    colframe=KazeInk,
    boxrule=1pt,
    sharp corners,
    fonttitle=\sffamily\bfseries\small,
    coltitle=white,
    attach boxed title to top left={yshift=-2mm, xshift=5mm},
    boxed title style={colback=KazeInk, sharp corners, boxrule=0pt},
    title=#1,
    left=5mm, top=6mm, bottom=4mm
}

\newcommand{\KazePageTitle}[1]{%
    \vspace{1em}
    {\sffamily\small\textcolor{KazeSecondary}{\textbf{\MakeUppercase{Document Scope: #1}}}\par}
    \vspace{0.5em}
}

\newcommand{\KazeDocumentNotice}{%
    \begin{kazecalloutbox}[Security Disclosure]
    \sffamily\small This dossier contains trade secrets and sensitive intellectual property. \textbf{Unauthorized distribution is strictly prohibited.} Access is limited to Kaze personnel and authorized external stakeholders for legal and investor review.
    \end{kazecalloutbox}
}

% --- THE "Blueprints" COVER PAGE ---
\newcommand{\KazeCover}[4]{
    \hypersetup{pageanchor=false}
    \begin{titlepage}
    \thispagestyle{empty}

% Background "Pillar" for institutional weight
    \begin{tikzpicture}[remember picture,overlay]
    \fill[KazePrimary] (current page.north west) rectangle ([xshift=0.5cm]current page.south west);
    \draw[KazePrimary!55, line width=1.8pt, opacity=0.22, line cap=round]
        ([xshift=0.08\paperwidth,yshift=-0.10\paperheight]current page.north west) --
        ([xshift=0.72\paperwidth,yshift=-0.10\paperheight]current page.north west);
    \draw[KazeSecondary!70, line width=1.1pt, opacity=0.19, line cap=round]
        ([xshift=0.12\paperwidth,yshift=-0.135\paperheight]current page.north west) --
        ([xshift=0.88\paperwidth,yshift=-0.135\paperheight]current page.north west);
    \draw[KazePrimary!55, line width=1.2pt, opacity=0.20]
        ([xshift=0.62\paperwidth,yshift=-0.06\paperheight]current page.north west) --
        ([xshift=0.76\paperwidth,yshift=-0.06\paperheight]current page.north west) --
        ([xshift=0.82\paperwidth,yshift=-0.11\paperheight]current page.north west) --
        ([xshift=0.93\paperwidth,yshift=-0.11\paperheight]current page.north west) --
        ([xshift=0.93\paperwidth,yshift=-0.18\paperheight]current page.north west) --
        ([xshift=0.82\paperwidth,yshift=-0.18\paperheight]current page.north west) --
        ([xshift=0.75\paperwidth,yshift=-0.24\paperheight]current page.north west) --
        ([xshift=0.58\paperwidth,yshift=-0.24\paperheight]current page.north west);
    \draw[KazeSecondary!75, line width=1.35pt, opacity=0.18, line cap=round]
        ([xshift=0.18\paperwidth,yshift=0.26\paperheight]current page.south west) --
        ([xshift=0.84\paperwidth,yshift=0.26\paperheight]current page.south west);
    \draw[KazePrimary!55, line width=1.8pt, opacity=0.20, line cap=round]
        ([xshift=0.14\paperwidth,yshift=0.21\paperheight]current page.south west) --
        ([xshift=0.62\paperwidth,yshift=0.21\paperheight]current page.south west);
    \end{tikzpicture}

    \vspace*{0.5cm}
    \begin{flushleft}
    \hspace{1cm}
    \begin{tabular}{@{}l l@{}}
    \raisebox{-0.3\height}{\includegraphics[width=2.2cm]{\KazeLogoPath}} &
    \begin{minipage}[b]{8cm}
    {\KazeWordmarkCover\par}
    \vspace{0.1cm}
    {\KazeByGaboCover\par}
    \end{minipage}
    \end{tabular}
    \end{flushleft}

    \vspace{3.5cm}

    \hspace{1cm}
    \begin{minipage}{0.85\textwidth}
    \raggedright
    {\KazeCoverLabel{OFFICIAL DOSSIER / #4}\par}
    \vspace{0.95cm}
    {\KazeCoverTitle{#1}\par}
    \vspace{0.7cm}
    {\KazeCoverSubtitle{#2}\par}

    \vspace{1.5cm}
    {\color{KazeLine}\hrule height 2pt width 3cm}
    \vspace{1.5cm}

    {\sffamily\large\textcolor{KazeMuted}{#3}\par}
    \end{minipage}

    \vfill

    \hspace{1cm}
    \begin{minipage}{0.85\textwidth}
    \begin{kazehighlightbox}
    \sffamily\small
    \textbf{\textcolor{KazeLegal}{RESTRICTED DOCUMENT:}} This information is intended for the recipient only. It contains proprietary strategic planning belonging to Kaze. By accessing this document, you agree to non-disclosure terms.
    \end{kazehighlightbox}
    \end{minipage}

    \vspace{1.5cm}
    \hspace{1cm}
    \begin{minipage}{0.85\textwidth}
    \sffamily\tiny\textcolor{KazeMuted}{AUTHORIZATION \& CONTACT:} \\
    \sffamily\small\textbf{\href{mailto:\KazeContactEmail}{\KazeContactEmail}} \hfill \textbf{\MakeUppercase{\KazeContactWebsiteLabel}}
    \end{minipage}

    \end{titlepage}
    \hypersetup{pageanchor=true}
}

\usepackage{listings}
\usepackage{upquote}

\definecolor{CodeBg}{HTML}{F4F1EA}
\definecolor{CodeRule}{HTML}{D8CCBC}
\definecolor{CodeComment}{HTML}{6A737D}
\definecolor{CodeKeyword}{HTML}{284F58}
\definecolor{CodeString}{HTML}{8A6538}

\lstdefinestyle{kazekotlin}{
    language={},
    basicstyle=\ttfamily\small,
    keywordstyle=\color{CodeKeyword}\bfseries,
    commentstyle=\color{CodeComment},
    stringstyle=\color{CodeString},
    showstringspaces=false,
    columns=fullflexible,
    keepspaces=true,
    breaklines=true,
    frame=single,
    rulecolor=\color{CodeRule},
    backgroundcolor=\color{CodeBg},
    xleftmargin=0.2cm,
    xrightmargin=0.2cm,
    aboveskip=0.8em,
    belowskip=0.8em
}

\begin{document}

\KazeCover
  {KMP Learning Guide for Kaze}
  {A course-style guide for building good Kotlin Multiplatform apps}
  {April 2026}
  {Developer learning workbook}

\KazePageTitle{KMP Learning Guide}
\KazeDocumentNotice

\section{How To Read This Guide}

This guide assumes you already know basic programming and have at least some Kotlin exposure. The goal here is not to explain what a class or function is. The goal is to teach how Kotlin Multiplatform apps are structured and why they are structured that way.

This document tries to answer five recurring questions:
\begin{itemize}
    \item what is this concept?
    \item why does KMP need it?
    \item how is it different from single-platform Android or iOS thinking?
    \item when should I use it?
    \item how does it show up in Kaze?
\end{itemize}

\begin{kazecalloutbox}[Important Mindset]
Kotlin Multiplatform is not magic code sharing. It is a design approach. You still need to decide what should be shared, what must stay platform-specific, and where UI, domain logic, and platform APIs should live.
\end{kazecalloutbox}

\section{Glossary}

\textbf{KMP}

Kotlin Multiplatform. A Kotlin approach for sharing some code across multiple platforms while still allowing platform-specific implementations.

\textbf{Compose Multiplatform}

JetBrains' multiplatform UI toolkit based on Compose concepts.

\textbf{Common Code}

Code in shared source sets such as \texttt{commonMain}, meant to work across platforms.

\textbf{Platform-Specific Code}

Code that only works on one target, such as Android, iOS, JVM, JS, or Wasm.

\textbf{Source Set}

A Gradle/Kotlin grouping of code for a target or shared layer, such as \texttt{commonMain}, \texttt{androidMain}, or \texttt{iosMain}.

\textbf{Composable}

A function annotated with \texttt{@Composable} that describes UI in Compose.

\textbf{State}

Data that affects what the UI shows.

\textbf{Recomposition}

Compose rerunning composable functions when relevant state changes.

\textbf{Expect/Actual}

Kotlin Multiplatform mechanism for declaring a common API and implementing it differently per platform.

\textbf{Platform Bridge}

A shared abstraction over a platform-specific capability, such as opening URLs, secure storage, haptics, or network reachability.

\textbf{Intent}

On Android, an object representing an action request, such as opening another app, sharing content, or launching a screen.

\textbf{Deep Link}

A link that opens a specific screen or flow inside an app.

\textbf{Canvas}

A low-level drawing surface used for custom graphics, maps, charts, and interactions.

\section{Chapter 1: What Kotlin Multiplatform Actually Is}

\subsection{The Core Idea}

Kotlin Multiplatform is a way to share code across platforms without pretending every platform is the same.

That is a very important point.

KMP does not say:
\begin{quote}
everything should be shared
\end{quote}

It says:
\begin{quote}
share what makes sense, and keep platform-specific things where they belong
\end{quote}

\subsection{What KMP Usually Shares Well}

Good candidates for shared code:
\begin{itemize}
    \item domain models
    \item business rules
    \item use cases
    \item repository interfaces
    \item networking models
    \item state management
    \item sometimes UI, if using Compose Multiplatform
\end{itemize}

\subsection{What Usually Stays Platform-Specific}

Things that often remain platform-specific:
\begin{itemize}
    \item Android intents
    \item iOS system integrations
    \item secure storage implementations
    \item push notifications
    \item camera and permissions details
    \item app entrypoints
\end{itemize}

\subsection{Why This Matters}

The biggest beginner mistake in KMP is trying to share code for ideological reasons instead of architectural reasons.

Good KMP design is not:
\begin{quote}
share as much as possible
\end{quote}

Good KMP design is:
\begin{quote}
share what improves maintainability without fighting the platforms
\end{quote}

\subsection{How Kaze Uses This}

Kaze already reflects this structure:
\begin{itemize}
    \item \texttt{shared} contains domain and shared logic
    \item \texttt{composeApp} contains Compose Multiplatform UI
    \item \texttt{iosApp} holds iOS entry code
    \item platform-specific code exists under \texttt{androidMain}, \texttt{iosMain}, \texttt{jsMain}, and \texttt{wasmJsMain}
\end{itemize}

\section{Chapter 2: Project Structure and Source Sets}

\subsection{Why Source Sets Exist}

KMP needs a way to organize code by sharing level.

That is what source sets do.

Common examples:
\begin{itemize}
    \item \texttt{commonMain}: shared production code
    \item \texttt{commonTest}: shared tests
    \item \texttt{androidMain}: Android-specific code
    \item \texttt{iosMain}: iOS-specific code
    \item \texttt{jsMain}: JavaScript-specific code
    \item \texttt{wasmJsMain}: WebAssembly-specific code
\end{itemize}

\subsection{How To Think About This}

Imagine concentric circles:
\begin{itemize}
    \item outer shared circle: code all platforms can use
    \item inner platform circles: code only one platform can use
\end{itemize}

If code imports Android-only APIs, it cannot live in \texttt{commonMain}.

\subsection{Best Practice}

\begin{itemize}
    \item keep domain and app logic as high in the shared hierarchy as reasonable
    \item move code down into platform source sets only when platform APIs force it
    \item do not pollute \texttt{commonMain} with fake abstractions for things that are deeply platform-native unless you truly need sharing
\end{itemize}

\subsection{A Tiny Source Set Example}

Shared code:

\begin{lstlisting}[style=kazekotlin]
// commonMain
data class UserProfile(
    val id: String,
    val displayName: String
)
\end{lstlisting}

Platform-specific code:

\begin{lstlisting}[style=kazekotlin]
// androidMain
fun defaultPlatformName(): String = "Android"

// iosMain
fun defaultPlatformName(): String = "iOS"
\end{lstlisting}

This is the basic KMP idea in miniature: share the parts that are truly common, split the parts that are platform-specific.

\section{Chapter 3: Compose Multiplatform and Composables}

\subsection{What A Composable Is}

A composable is a function that describes UI.

\begin{lstlisting}[style=kazekotlin]
@Composable
fun Greeting(name: String) {
    Text("Hello, $name")
}
\end{lstlisting}

The important mindset shift is this:

Compose UI is not about manually mutating views step by step.
It is about describing what the UI should look like for the current state.

\subsection{A Slightly More Real Example}

\begin{lstlisting}[style=kazekotlin]
@Composable
fun ProfileCard(
    name: String,
    isOnline: Boolean,
    onMessageClick: () -> Unit
) {
    Column {
        Text(text = name)
        Text(
            text = if (isOnline) "Online" else "Offline"
        )
        Button(onClick = onMessageClick) {
            Text("Message")
        }
    }
}
\end{lstlisting}

What this example teaches:
\begin{itemize}
    \item composables receive data as parameters
    \item composables receive actions as callbacks
    \item the UI is derived from current state such as \texttt{isOnline}
\end{itemize}

\subsection{Why This Style Exists}

Traditional UI systems often involve:
\begin{itemize}
    \item inflating layouts
    \item finding views
    \item mutating them manually
    \item keeping view state and app state synchronized
\end{itemize}

Compose tries to reduce that complexity by making the UI a function of state.

\subsection{What Recomposition Means}

When state changes, Compose may run relevant composables again.

This is called recomposition.

Beginners often worry:
\begin{quote}
if functions rerun, is that bad?
\end{quote}

No. Recomposition is normal. Compose is designed around it.

What matters is keeping composables predictable and making sure heavy work is not repeated carelessly.

\subsection{What To Watch Out For}

\begin{itemize}
    \item do not put heavy blocking work directly in composables
    \item keep composables as UI descriptions, not giant business-logic containers
    \item remember that local state and remembered objects affect recomposition behavior
\end{itemize}

\subsection{How Kaze Uses This}

Kaze's main app entry in [KazeApp.kt](/Users/muhirwagabooreste/AndroidStudioProjects/kaze/composeApp/src/commonMain/kotlin/dev/orestegabo/kaze/KazeApp.kt) shows a large Compose tree with shared screens, state-driven navigation, and view models.

\section{Chapter 4: State, Recomposition, and View Models}

\subsection{What State Is}

State is any data that changes what the UI shows.

Examples:
\begin{itemize}
    \item whether a user is signed in
    \item which tab is selected
    \item whether a network call is loading
    \item which map room is selected
\end{itemize}

\subsection{Local State Example}

\begin{lstlisting}[style=kazekotlin]
@Composable
fun CounterCard() {
    var count by remember { mutableStateOf(0) }

    Column {
        Text("Count: $count")
        Button(onClick = { count++ }) {
            Text("Increase")
        }
    }
}
\end{lstlisting}

This works because:
\begin{itemize}
    \item \texttt{count} is state
    \item Compose tracks it
    \item when it changes, the relevant UI recomposes
\end{itemize}

\subsection{Screen State Example}

\begin{lstlisting}[style=kazekotlin]
data class ProfileUiState(
    val isLoading: Boolean = true,
    val displayName: String = "",
    val errorMessage: String? = null
)
\end{lstlisting}

This kind of state belongs at the screen or view-model level because it represents the whole screen's condition, not just one tiny widget.

\subsection{Why State Management Is Central}

Most UI bugs are really state bugs:
\begin{itemize}
    \item state changed but UI did not update
    \item UI updated from stale state
    \item too much mutable state exists in too many places
\end{itemize}

Compose encourages you to make state explicit.

\subsection{Local State vs Screen State}

Local state is good for small UI concerns:
\begin{itemize}
    \item expanded/collapsed flag
    \item text field content in a local form
\end{itemize}

Screen or app state often belongs in a view model or similar controller:
\begin{itemize}
    \item authenticated session
    \item loaded domain data
    \item navigation-related app mode
\end{itemize}

\subsection{Why View Models Help}

View models help separate:
\begin{itemize}
    \item UI description
    \item state ownership
    \item business logic coordination
\end{itemize}

This is especially valuable in larger apps because composables stay focused on rendering.

\subsection{A View Model Style Example}

\begin{lstlisting}[style=kazekotlin]
class ProfileViewModel(
    private val repository: ProfileRepository
) {
    var uiState by mutableStateOf(ProfileUiState())
        private set

    suspend fun loadProfile() {
        uiState = uiState.copy(isLoading = true)
        uiState = try {
            val profile = repository.loadProfile()
            uiState.copy(
                isLoading = false,
                displayName = profile.displayName,
                errorMessage = null
            )
        } catch (e: Exception) {
            uiState.copy(
                isLoading = false,
                errorMessage = "Could not load profile"
            )
        }
    }
}
\end{lstlisting}

This example is useful because it shows the split clearly:
\begin{itemize}
    \item repository loads data
    \item view model manages screen state
    \item composable reads and renders the state
\end{itemize}

\subsection{How Kaze Uses This}

Kaze has view-model-style classes such as \texttt{StayViewModel}, \texttt{EventsViewModel}, \texttt{ExploreViewModel}, and \texttt{KazeAppViewModel}. That is a healthy sign that the app is not trying to put all behavior inside composables.

\section{Chapter 5: Navigation in KMP Apps}

\subsection{Why Navigation Matters}

Navigation is not just changing screens. It is managing app flow.

Good navigation architecture affects:
\begin{itemize}
    \item how screens communicate
    \item how deep links work
    \item how state survives transitions
    \item how easy the app is to test
\end{itemize}

\subsection{Shared Navigation vs Platform Navigation}

In a KMP app, you can:
\begin{itemize}
    \item keep navigation mostly shared in common code
    \item delegate some navigation concerns to platform-specific entrypoints
\end{itemize}

This depends on your architecture.

\subsection{A Simple Shared Navigator Pattern}

Kaze uses a shared navigator-style object in [KazeNavigator.kt](/Users/muhirwagabooreste/AndroidStudioProjects/kaze/composeApp/src/commonMain/kotlin/dev/orestegabo/kaze/presentation/navigation/KazeNavigator.kt).

That approach works well when:
\begin{itemize}
    \item the app's destinations are known in shared code
    \item UI flow is mostly common across platforms
    \item you want state-driven navigation rather than platform-fragment-driven navigation
\end{itemize}

\subsection{Simple Navigation Example}

\begin{lstlisting}[style=kazekotlin]
enum class AppDestination {
    HOME,
    PROFILE,
    SETTINGS
}

class AppNavigator {
    var currentDestination by mutableStateOf(AppDestination.HOME)
        private set

    fun goTo(destination: AppDestination) {
        currentDestination = destination
    }
}
\end{lstlisting}

In the UI:

\begin{lstlisting}[style=kazekotlin]
@Composable
fun AppRoot(navigator: AppNavigator) {
    when (navigator.currentDestination) {
        AppDestination.HOME -> HomeScreen(
            onOpenProfile = { navigator.goTo(AppDestination.PROFILE) }
        )
        AppDestination.PROFILE -> ProfileScreen()
        AppDestination.SETTINGS -> SettingsScreen()
    }
}
\end{lstlisting}

This is not the only navigation style in KMP, but it is one of the easiest to understand.

\subsection{Best Practice}

\begin{itemize}
    \item model destinations clearly
    \item avoid letting every screen navigate however it wants with hidden global side effects
    \item keep navigation state understandable and testable
\end{itemize}

\section{Chapter 6: Platform Actions, Android Intents, and Their KMP Equivalents}

\subsection{What An Android Intent Is}

On Android, an intent is a message-like object used to request an action:
\begin{itemize}
    \item open another app
    \item share content
    \item open a browser
    \item launch another activity
\end{itemize}

\subsection{Why This Becomes A KMP Question}

In a multiplatform app, Android intents are not directly available in common code.
So the real question becomes:

\begin{quote}
what is the shared abstraction for ``ask the platform to do something external''?
\end{quote}

\subsection{The KMP Way To Think About It}

Instead of putting Android intent code directly in shared UI logic, you usually create a shared abstraction such as:
\begin{itemize}
    \item open external URL
    \item share text
    \item launch auth flow
    \item open map app
\end{itemize}

Then each platform implements that behavior in its own way.

\subsection{What These Are Called Outside Android}

They are not usually called ``intents'' on other platforms.

Instead, you might think in terms of:
\begin{itemize}
    \item URL launching
    \item deep links
    \item navigation requests
    \item platform service actions
    \item OS integration APIs
\end{itemize}

\subsection{Expect/Actual and Platform Bridges}

One strong KMP pattern is:
\begin{enumerate}
    \item define a shared capability
    \item provide platform-specific implementations
\end{enumerate}

Example concept:

\begin{lstlisting}[style=kazekotlin]
interface ExternalUrlLauncher {
    fun open(url: String): Boolean
}
\end{lstlisting}

This is usually a better shared design than exposing Android intent details directly to common code.

\subsection{A Shared Platform Action Example}

\begin{lstlisting}[style=kazekotlin]
@Composable
fun HelpButton(urlLauncher: ExternalUrlLauncher) {
    Button(onClick = {
        urlLauncher.open("https://kazerwanda.com/help")
    }) {
        Text("Open help")
    }
}
\end{lstlisting}

This is a good KMP example because the shared UI says what it wants:
\begin{quote}
open this help link
\end{quote}

but it does not care whether Android uses an intent, iOS uses its own system API, or web uses browser logic.

\subsection{How Kaze Uses This}

Kaze already uses platform-bridge ideas around auth and external launching. That is exactly the sort of thing you want in KMP: shared app logic asks for a capability, while platform-specific code handles the real OS integration.

\subsection{Best Practice}

\begin{itemize}
    \item expose shared intent-like behavior as capabilities, not as Android-specific types
    \item keep the Android intent object itself in Android code
    \item keep common code focused on the action the app wants, not the platform mechanism
\end{itemize}

\section{Chapter 7: Resources, Themes, and Design Systems}

\subsection{Why This Matters}

A good KMP app is not just shared code. It is a coherent cross-platform product.

That means you need a way to think about:
\begin{itemize}
    \item colors
    \item typography
    \item spacing
    \item icons
    \item images
    \item strings
\end{itemize}

\subsection{Shared Design Language}

Compose Multiplatform makes it natural to share a design system:
\begin{itemize}
    \item theme definitions
    \item common components
    \item layout primitives
    \item visual states
\end{itemize}

\subsection{How Kaze Uses This}

Kaze's shared UI and theme files show this direction clearly: common theme definitions, shared screens, and reusable UI components across the app.

\section{Chapter 8: Networking, Repositories, and App Architecture}

\subsection{Why Architecture Matters More In KMP}

Because KMP apps already have shared and platform-specific layers, architecture becomes even more important.

Without clear boundaries, your code can become confusing quickly.

\subsection{A Healthy Layering Pattern}

Many good KMP apps separate:
\begin{itemize}
    \item UI
    \item presentation/state management
    \item domain/use cases
    \item data/repositories
    \item platform services
\end{itemize}

\subsection{Why Repositories Matter}

Repositories let shared business logic depend on abstract data access rather than platform-specific implementation details.

That means:
\begin{itemize}
    \item UI stays cleaner
    \item fake repositories are easier in tests
    \item storage and networking decisions stay replaceable
\end{itemize}

\section{Chapter 9: Platform Services and Expect/Actual}

\subsection{What Expect/Actual Solves}

Sometimes a shared module needs a capability, but the implementation must differ by platform.

That is what \texttt{expect}/\texttt{actual} is for.

\subsection{Mental Model}

\begin{itemize}
    \item \texttt{expect}: common code declares what it needs
    \item \texttt{actual}: each platform provides the real implementation
\end{itemize}

\subsection{When To Use It}

Use \texttt{expect}/\texttt{actual} when:
\begin{itemize}
    \item the capability is truly shared in concept
    \item the implementation must differ by platform
\end{itemize}

Examples:
\begin{itemize}
    \item platform name or device label
    \item secure store implementation
    \item network reachability
    \item browser or external URL launching
\end{itemize}

\subsection{Expect/Actual Example}

\begin{lstlisting}[style=kazekotlin]
// commonMain
expect fun defaultDeviceLabel(): String
\end{lstlisting}

\begin{lstlisting}[style=kazekotlin]
// androidMain
actual fun defaultDeviceLabel(): String = "Android device"
\end{lstlisting}

\begin{lstlisting}[style=kazekotlin]
// iosMain
actual fun defaultDeviceLabel(): String = "iPhone"
\end{lstlisting}

This is one of the clearest examples of what KMP is trying to do:
\begin{itemize}
    \item common code defines the need
    \item each platform supplies the real answer
\end{itemize}

\subsection{When Not To Force It}

Do not create an \texttt{expect}/\texttt{actual} abstraction for every tiny platform difference if a simpler architecture would be clearer.

\section{Chapter 10: Canvas, Custom Drawing, and Interactive Graphics}

\subsection{What Canvas Is}

Canvas is a low-level drawing API. Instead of composing only standard UI components like rows, columns, buttons, and text fields, you draw shapes, lines, paths, gradients, and labels manually.

\subsection{A Tiny Canvas Example}

\begin{lstlisting}[style=kazekotlin]
@Composable
fun SimpleBadge() {
    Canvas(modifier = Modifier.size(120.dp)) {
        drawCircle(
            color = Color(0xFF284F58),
            radius = size.minDimension / 2
        )
        drawLine(
            color = Color.White,
            start = Offset(20f, size.height / 2),
            end = Offset(size.width - 20f, size.height / 2),
            strokeWidth = 8f
        )
    }
}
\end{lstlisting}

This example is intentionally small. It teaches the core idea:
\begin{itemize}
    \item you are drawing directly
    \item you work with coordinates and sizes
    \item there are no ready-made widgets doing the visual work for you
\end{itemize}

\subsection{Why You Would Use Canvas}

Use canvas when the UI is:
\begin{itemize}
    \item map-like
    \item chart-like
    \item diagram-like
    \item game-like
    \item highly custom
\end{itemize}

\subsection{Why Canvas Feels Different From Normal Compose UI}

In normal Compose UI, you describe ready-made widgets.

In canvas work, you describe geometry and drawing operations.

So you need to think about:
\begin{itemize}
    \item coordinates
    \item bounds
    \item scaling
    \item transforms
    \item hit-testing
    \item custom gestures
\end{itemize}

\subsection{A Map-Like Interaction Example}

\begin{lstlisting}[style=kazekotlin]
var zoom by remember { mutableFloatStateOf(1f) }
var pan by remember { mutableStateOf(Offset.Zero) }

Canvas(
    modifier = Modifier.pointerInput(Unit) {
        detectTransformGestures { _, panChange, zoomChange, _ ->
            zoom *= zoomChange
            pan += panChange
        }
    }
) {
    withTransform({
        translate(pan.x, pan.y)
        scale(zoom, zoom)
    }) {
        // draw shapes here
    }
}
\end{lstlisting}

This is the kind of pattern you use for zoomable maps, diagrams, and floor plans.

\subsection{How Kaze Uses This}

Kaze's [KazeMapView.kt](/Users/muhirwagabooreste/AndroidStudioProjects/kaze/composeApp/src/commonMain/kotlin/dev/orestegabo/kaze/ui/map/components/KazeMapView.kt) is a strong example of a KMP-friendly canvas use case: shared custom drawing for map-style interaction, including zoom, pan, selection, and custom rendering.

\subsection{What To Pay Attention To}

\begin{itemize}
    \item keep drawing math well-structured
    \item separate rendering logic from domain data when possible
    \item be careful with performance and repeated allocations
    \item think about gesture input and coordinate conversion clearly
\end{itemize}

\section{Chapter 11: Gestures, Animations, and Interaction}

\subsection{Why Interaction Design Matters}

A good KMP app is not only correct. It must also feel responsive and natural.

This includes:
\begin{itemize}
    \item taps
    \item scroll
    \item drag
    \item pinch-to-zoom
    \item transitions
    \item loading feedback
\end{itemize}

\subsection{What To Watch Out For}

\begin{itemize}
    \item gesture code can become hard to reason about if mixed carelessly with business state
    \item animations should support meaning, not just decoration
    \item custom interactions should still feel usable across platforms
\end{itemize}

\section{Chapter 12: Lifecycle, Side Effects, and Asynchronous Work}

\subsection{Why Side Effects Need Care}

UI code is not just rendering. Sometimes it must:
\begin{itemize}
    \item load data
    \item start a network request
    \item react to external events
    \item persist something
\end{itemize}

These are side effects.

\subsection{Why This Is Tricky In Compose}

Because composables may re-run, you must think carefully about when work should happen.

That is why Compose has effect APIs and why state ownership matters so much.

\subsection{Coroutines In KMP Apps}

Coroutines are a major part of KMP app development because they help coordinate:
\begin{itemize}
    \item network calls
    \item storage access
    \item background work
    \item cancellation-aware UI actions
\end{itemize}

\subsection{A Side Effect Example}

\begin{lstlisting}[style=kazekotlin]
@Composable
fun ProfileScreen(viewModel: ProfileViewModel) {
    LaunchedEffect(Unit) {
        viewModel.loadProfile()
    }

    Text(viewModel.uiState.displayName)
}
\end{lstlisting}

This example teaches an important Compose idea:
\begin{itemize}
    \item composables describe UI
    \item side effects such as loading data should happen through effect APIs
    \item \texttt{LaunchedEffect} gives you a safe place to start coroutine-based work tied to composition
\end{itemize}

\section{Chapter 13: Testing KMP Apps}

\subsection{What To Test}

A good KMP app usually needs tests at multiple levels:
\begin{itemize}
    \item shared logic tests
    \item view-model or state tests
    \item repository tests
    \item UI tests where practical
\end{itemize}

\subsection{A Shared Logic Test Example}

\begin{lstlisting}[style=kazekotlin]
class GreetingFormatterTest {
    @Test
    fun formatsName() {
        val result = formatGreeting("Aline")
        assertEquals("Hello, Aline", result)
    }
}
\end{lstlisting}

This may look simple, but it demonstrates one of KMP's big strengths:
\begin{quote}
logic in shared code can be tested once and reused across targets
\end{quote}

\subsection{Why Shared Logic Testing Is Powerful}

One of KMP's strengths is that shared logic can be tested once and benefit multiple targets.

That makes good architectural separation even more valuable.

\section{Chapter 14: Performance, Debugging, and Common Mistakes}

\subsection{Common Mistakes}

\begin{longtable}{p{0.26\textwidth} p{0.28\textwidth} p{0.36\textwidth}}
\textbf{Mistake} & \textbf{Symptom} & \textbf{What To Watch} \\
Trying to share everything & Architecture feels forced & Accept platform differences where appropriate. \\
Putting too much logic in composables & Hard-to-maintain UI & Move state and logic into view models or use cases. \\
Platform APIs leaking into common code & Build/design problems & Use abstractions or platform source sets. \\
Bad canvas/state separation & Messy interactive graphics & Separate geometry, rendering, and interaction concerns. \\
Ignoring recomposition behavior & Performance or correctness surprises & Understand state and effect APIs. \\
\end{longtable}

\subsection{Debugging Habit}

When a KMP UI problem appears, ask:
\begin{enumerate}
    \item is the state correct?
    \item is the state in the right place?
    \item is recomposition happening as expected?
    \item is this actually a platform-specific bug rather than a shared-code bug?
\end{enumerate}

\section{Chapter 15: What You Need To Build A Good KMP App}

\subsection{Core Ingredients}

A good KMP app usually needs:
\begin{itemize}
    \item clear shared domain models
    \item strong state management
    \item composables that stay mostly declarative
    \item a navigation strategy
    \item platform bridges for OS-specific capabilities
    \item repository-driven data flow
    \item good coroutine usage
    \item testing strategy
    \item design-system consistency
\end{itemize}

\subsection{What Makes KMP Apps Go Wrong}

KMP apps often go wrong when teams:
\begin{itemize}
    \item confuse code sharing with good architecture
    \item let platform details leak everywhere
    \item make UI composables do too much
    \item ignore how state and side effects interact
\end{itemize}

\section{Mapping This Back To Kaze}

\begin{longtable}{p{0.28\textwidth} p{0.64\textwidth}}
\textbf{Concept} & \textbf{Where To Look In Kaze} \\
Shared UI app root & \texttt{composeApp/src/commonMain/kotlin/dev/orestegabo/kaze/KazeApp.kt} \\
Shared navigation & \texttt{composeApp/src/commonMain/kotlin/dev/orestegabo/kaze/presentation/navigation/KazeNavigator.kt} \\
Shared dependencies assembly & \texttt{composeApp/src/commonMain/kotlin/dev/orestegabo/kaze/presentation/di/KazeDependencies.kt} \\
Platform-specific implementations & \texttt{androidMain}, \texttt{iosMain}, \texttt{jsMain}, \texttt{wasmJsMain} source sets \\
Canvas-based rendering & \texttt{composeApp/src/commonMain/kotlin/dev/orestegabo/kaze/ui/map/components/KazeMapView.kt} \\
Shared domain layer & \texttt{shared/src/commonMain/kotlin/dev/orestegabo/kaze/...} \\
\end{longtable}

\section{Exercises}

\subsection{Architecture}

\begin{itemize}
    \item List three things in an app that should probably live in shared code.
    \item List three things that should probably stay platform-specific.
\end{itemize}

\subsection{Compose}

\begin{itemize}
    \item Write a simple composable that shows a title and button.
    \item Explain in your own words what recomposition means.
\end{itemize}

\subsection{Platform Bridges}

\begin{itemize}
    \item Design a shared interface for ``open external URL'' without mentioning Android intents directly.
    \item Explain why that is a better shared abstraction than exposing Android APIs to common code.
\end{itemize}

\subsection{Canvas}

\begin{itemize}
    \item Name two app features that would be better built with canvas than with only standard widgets.
    \item Explain why zoomable maps are a strong canvas use case.
\end{itemize}

\subsection{State and Navigation}

\begin{itemize}
    \item Explain the difference between local UI state and screen-level app state.
    \item Explain one advantage of having a shared navigator object.
\end{itemize}

\section{Final Advice}

The healthiest way to learn KMP is to stop asking:
\begin{quote}
how much code can I share?
\end{quote}

and start asking:
\begin{quote}
what architecture gives me the clearest, safest, and most maintainable cross-platform app?
\end{quote}

That is the mindset that leads to good KMP applications.

\section{Contact}

\KazeContactBlock

\end{document}
